\chapter{Attacks on Networks, with Focus on DDoS}
\label{ch:attacks-ddos}

The previous chapter described what a detector observes. This chapter
describes what it has to find. It fixes the attack families used as labels,
looks closely at denial of service, which dominates the incident figures of
\cref{ch:introduction}, and covers scanning, credential attacks and privilege
escalation. It then turns to the defences and ends with the domain knowledge
that can be written as logical rules.

% ================================================================
% 3.1 Attack taxonomy used in this thesis
% ================================================================
\section{Attack taxonomy used in this thesis}
\label{sec:attack-taxonomy}

A labelled dataset asks the detector to separate a fixed set of classes, so
the taxonomy behind the labels decides what is being tested. This thesis uses
the four families of KDD and NSL-KDD as its base taxonomy, because the labels
of NSL-KDD are defined this way. Each family is then given a name in MITRE
ATT\&CK, so that later chapters can refer to attacks in a shared vocabulary.
The two taxonomies are built on different principles. This section shows
where they agree and where they do not.

\subsection{The classic KDD families: DoS, Probe, R2L, U2R}
\label{sec:kdd-families}

The task description of the KDD Cup 1999 groups the attacks in four
categories \citep{kddcup1999task}. Denial of service is the first, and SYN
flood is the example it gives. Probing is described as surveillance and other
probing, with port scanning as the example. Remote-to-Local, or R2L, is
unauthorized access from a remote machine, and the example is password
guessing. User-to-Root, or U2R, is unauthorized access to local superuser
privileges, and the example is a buffer overflow.

The same document separates the families by how they appear in the traffic.
DoS and probing attacks involve many connections to a host in a short period.
R2L and U2R attacks do not. They are embedded in the data portion of the
packets and normally involve a single connection
\citep{kddcup1999task,tavallaee2009kddcup99}. For a detector limited to
flow- and header-level features, this split has a direct consequence. Counts
and rates carry the trace of the first group. For the second group the
evidence sits in the content, which these detectors do not see. This is our
reading, and \cref{sec:r2l-u2r} comes back to it.

The KDD family is called DoS. The rest of this thesis reserves DoS for
attacks from a single source and DDoS for attacks from many sources. When the
whole KDD family is meant, the text says so.

\subsection{Mapping the families onto MITRE ATT\&CK tactics}
\label{sec:families-attack-mapping}

ATT\&CK describes adversary behaviour by tactic, the tactical objective of
the adversary, and by technique, the way the objective is reached
\citep{strom2020attack}. The Enterprise matrix lists the tactics
\citep{mitre2026tactics}. Six of them matter here. Reconnaissance gathers
information to plan future operations. Initial Access is the attempt to get
into the network. Credential Access is the theft of account names and
passwords. Discovery is the attempt to figure out the environment. Privilege
Escalation seeks higher-level permissions. Impact manipulates, interrupts or
destroys systems and data.

\Cref{tab:kdd-attack-mapping} aligns each family with these tactics and with
the techniques that fit its KDD example. Denial of service falls under
Impact. Technique T1498 covers network floods, and T1499 covers attacks on
the services of an endpoint
\citep{mitre2025t1498,mitre2025t1499}. Probing is split in two. Active
scanning of victim infrastructure is Reconnaissance, technique T1595
\citep{mitre2025t1595}. Network service discovery is Discovery, technique
T1046, and it lists port scans among its methods \citep{mitre2026t1046}.
Password guessing is Credential Access, as a sub-technique of brute force,
T1110 \citep{mitre2026t1110}. Exploiting a weakness in an Internet-facing
system is Initial Access, T1190 \citep{mitre2026t1190}. A buffer overflow that
gives root fits Privilege Escalation, T1068 \citep{mitre2025t1068}.

\begin{table}[t]
\centering
\small
\caption{KDD attack families aligned with ATT\&CK tactics and techniques.
The alignment is ours.}
\label{tab:kdd-attack-mapping}
\begin{tabular}{@{}p{1.3cm}p{2.7cm}p{3cm}p{5.2cm}@{}}
\toprule
Family & KDD example & ATT\&CK tactic & ATT\&CK technique \\
\midrule
DoS   & SYN flood & Impact
      & T1498 Network Denial of Service; T1499 Endpoint Denial of Service \\
Probe & Port scanning & Reconnaissance; Discovery
      & T1595 Active Scanning; T1046 Network Service Discovery \\
R2L   & Password guessing & Credential Access; Initial Access
      & T1110 Brute Force; T1190 Exploit Public-Facing Application \\
U2R   & Buffer overflow & Privilege Escalation
      & T1068 Exploitation for Privilege Escalation \\
\bottomrule
\end{tabular}
\end{table}

The alignment is not one to one. The names R2L and U2R carry the starting
point of the attack, a remote machine or a local user. ATT\&CK tactics name
the objective. Two families therefore map to two tactics each. A scan is
Reconnaissance when it comes from outside and Discovery when the adversary
already holds a host. A detector that sees only the traffic of the scan may
not be able to tell the two apart. Neither KDD nor ATT\&CK states this
alignment, and it is ours. ATT\&CK serves here as a label space, as in the
previous chapter. The conditions a network detector can actually test come
from the next sections, and \cref{sec:domain-knowledge-formalisation} turns
them into candidate axioms.

% ================================================================
% 3.2 Denial of Service and Distributed Denial of Service
% ================================================================
\section{Denial of Service and Distributed Denial of Service}
\label{sec:dos-ddos}

A denial of service attack tries to keep legitimate users away from a
service. To reach a large scale, attackers compromise many hosts and use them
together, and the attack is then distributed \citep{handley2006rfc4732}. The
two cases keep the names fixed in the previous section, DoS for one source
and DDoS for many. The attacks of this section differ above all in what they
exhaust: link capacity, protocol state or application resources. The
subsections follow this grouping, which is ours.

\subsection{Volumetric attacks: UDP flood, ICMP flood, amplification/\allowbreak reflection (DNS, NTP, memcached, LDAP, SSDP)}
\label{sec:volumetric-attacks}

The simplest attack sends enough traffic that does not slow down under
congestion to saturate a link \citep{handley2006rfc4732}. Mirkovic and Reiher
list the UDP flood and the ICMP echo flood as attacks that send packets for
non-critical services of the victim \citep{mirkovic2004taxonomy}. They call an
attack brute force when it initiates a vast amount of seemingly legitimate
transactions. The victim can do little on its own against such an attack,
because filtering also hits legitimate traffic and its resources are limited
\citep{mirkovic2004taxonomy}.

Reflection and amplification change both who sends the traffic and how much
of it there is. The attacker sends requests to third-party servers with the
address of the victim as source, and the servers reply to the victim
\citep{rossow2014amplification}. Rossow measures the bandwidth amplification
factor, the ratio between the UDP payload of a reply and that of the request,
for 14 protocols. In his results NTP reaches 556.9 over all the amplifiers
found and SSDP 30.8. Memcached went further. In the attack on GitHub of 28
February 2018 the factor was up to 51{,}000 and the traffic peaked at
1.35~Tbps \citep{kottler2018github}. DNS, LDAP, NTP and SSDP are among the
reflection attacks of the CIC-DDoS2019 dataset \citep{cic2019ddos}.

\subsection{Protocol / state-exhaustion attacks: SYN flood, ACK flood, fragmentation (teardrop), Slowloris and low-and-slow variants}
\label{sec:state-exhaustion-attacks}

A SYN flood does not try to fill the link. It exhausts the backlog of
half-open connections of a port. The attacker sends a barrage of SYN
segments, often with spoofed source addresses, that never answer the SYN-ACK.
Legitimate connection attempts are then rejected \citep{eddy2007rfc4987}.
Mirkovic and Reiher classify it as a semantic attack, since it exploits the
way TCP allocates space for connection requests
\citep{mirkovic2004taxonomy}. The attack commands of Mirai recorded by
Antonakakis et al. include ACK floods, ACK-STOMP floods and SYN floods
\citep{antonakakis2017mirai}.

Fragmentation attacks target reassembly. In the teardrop attack, some
implementations of the IP fragment reassembly code did not handle
overlapping fragments properly, and any remote user could crash a vulnerable
machine \citep{cert1997teardrop}.
Slowloris and the low-and-slow variants need little bandwidth. Cambiaso et
al. describe slow DoS attacks as exploiting application-layer protocols and
bringing down a well-equipped server with a small bandwidth
\citep{cambiaso2013slow}. Like a SYN flood, they work on the resources of the
server and not on the capacity of the link. This inference is ours.

\subsection{Application-layer attacks: HTTP flood, HTTP/2 rapid reset style attacks}
\label{sec:application-layer-attacks}

An HTTP flood sends requests that look like normal use of a web service.
Mirkovic and Reiher call such attacks non-filterable, because their packets
request legitimate services from the victim and cannot be told from those of
a legitimate client. They also note that HTTP flood packets must belong to an
established TCP connection and therefore cannot carry spoofed source
addresses \citep{mirkovic2004taxonomy}.
Booters list HTTP-based attacks, Slowloris and RUDY among their
application-layer attacks \citep{santanna2015inside}.

HTTP/2 rapid reset shows how an application protocol can become the weapon.
An HTTP/2 connection can carry many concurrent streams, and a stream can be
ended at once with a reset frame \citep{thomson2022rfc9113}. The weakness
lets an attacker reset many streams quickly, and the server consumes
resources in response. It was exploited in the wild from August to October
2023 \citep{nvd2023rapidreset}.

\subsection{Botnets, booters and DDoS-as-a-service}
\label{sec:botnets-booters}

Large attacks need many sources, and botnets supply them. At its peak Mirai
counted about 600{,}000 infected devices, mostly embedded and IoT devices, and
used them to overwhelm several high-profile targets
\citep{antonakakis2017mirai}. Booters make such attacks available to anyone.
They are websites that offer DDoS attacks as a paid service, with prices
starting at one dollar, and they remove the need for technical skills
\citep{santanna2015inside}.

\subsection{Observable signatures at flow level (the features that actually move)}
\label{sec:flow-signatures}

The attacks above act on different resources, but a flow-based detector sees
only their traces. \Cref{tab:ddos-flow-signatures} lists what changes in the
flow records of each. The mechanisms come from the cited sources, and the
traces are our interpretation of them.

\begin{table}[t]
\centering
\small
\caption{Traces of the attacks of this section in flow records. The traces
are derived by us from the cited mechanisms.}
\label{tab:ddos-flow-signatures}
\begin{tabular}{@{}p{2.8cm}p{6.5cm}p{2.9cm}@{}}
\toprule
Attack & What moves in the flow records & Mechanism \\
\midrule
UDP and ICMP flood & Very high packet and byte rate towards one destination,
  often from many sources
  & \citep{handley2006rfc4732,mirkovic2004taxonomy} \\
Reflection and amplification & Many flows from different source addresses,
  from the service port of one protocol, with large replies to requests that
  the destination never sent
  & \citep{rossow2014amplification} \\
SYN flood & Many short flows of SYN packets that never complete, more SYN than
  FIN and RST, often many distinct source addresses
  & \citep{eddy2007rfc4987,wang2002detecting} \\
Slow DoS & Low bit rate and small attacker bandwidth; long-lived flows (our inference)
  & \citep{cambiaso2013slow} \\
HTTP flood & Valid source addresses and complete connections, with a high
  request rate
  & \citep{mirkovic2004taxonomy} \\
HTTP/2 rapid reset & Few connections that carry very many requests and
  resets, which a flow record does not show
  & \citep{thomson2022rfc9113,nvd2023rapidreset} \\
\bottomrule
\end{tabular}
\end{table}

Two patterns stand out. Volumetric and state-exhaustion attacks move counts
and rates, so a threshold already catches part of the traffic. Reflection and
spoofed SYN floods also change the number of distinct source addresses, which
is why the address distributions of \cref{sec:statistical-anomaly-detection}
are informative. Slow and application-layer attacks can leave the volume
almost unchanged. They differ from normal traffic in the behaviour of each
source or connection, which a flow record shows only in part. These are the
conditions that the axioms of \cref{sec:domain-knowledge-formalisation}
formalise.

% ================================================================
% 3.3 Reconnaissance and scanning (Probe)
% ================================================================
\section{Reconnaissance and scanning (Probe)}
\label{sec:recon-scanning}

\subsection{Horizontal, vertical and stealth scanning; nmap/\allowbreak satan/\allowbreak ipsweep patterns}
\label{sec:scan-types}

A scan maps the target before other attacks. The review of Pittman reports
that Barnett et al. add vertical and horizontal scans to the classification of
De Vivo et al. \citep{pittman2023portscans}. This thesis calls a scan horizontal when it
probes the same port on many hosts and vertical when it probes many ports of
one host.

A scan can be stealthy in two ways. One concerns the packets. Pittman reports
a classification that sets TCP connect and SYN scans as classical against FIN,
XMAS and NULL scans, which it calls stealth \citep{pittman2023portscans}. The
other concerns time. A scanner can space its probes so that no short window
shows a pattern, which is why the KDD features add a window of 100
connections \citep{tavallaee2009kddcup99}.

The documentation of the DARPA evaluation behind the KDD data lists ipsweep,
nmap and satan among the probes \citep{kddcup1999task,mitll1998attacks}. An
ipsweep is a surveillance sweep to find the hosts that are listening on a
network \citep{mitll1998attacks}. We judge it close to a horizontal scan. Several of the studies reviewed by Pittman generate their scan
traffic with Nmap \citep{pittman2023portscans}.

At flow level we expect a scan to be a burst of short flows towards many
addresses or ports. This agrees with Lakhina et al., in whose data
the port scans and network scans were found only with entropy and not with
volume metrics \citep{lakhina2005mining}.

% ================================================================
% 3.4 Remote-to-Local and User-to-Root
% ================================================================
\section{Remote-to-Local and User-to-Root}
\label{sec:r2l-u2r}

\subsection{Credential attacks, guessed passwords, FTP/\allowbreak IMAP exploitation}
\label{sec:credential-attacks}

R2L attacks give a remote adversary access to a local account. Password
guessing, the KDD example, repeats attempts until one works, and ATT\&CK
files it under brute force \citep{mitre2026t1110}. The other path is a flaw in
an Internet-facing system. ATT\&CK lists web servers, databases and standard
services such as SMB and SSH among the exploited applications
\citep{mitre2026t1190}. The documentation of the DARPA evaluation gives
examples of both paths. A dictionary attack makes repeated guesses at user
names and passwords. Ftp-write takes advantage of a common misconfiguration
of anonymous FTP, and imap exploits a buffer overflow in an IMAP server
\citep{mitll1998attacks}. In our reading the two paths leave different traces.
Guessing repeats attempts and raises counts, while an exploit fits in one
connection and does not.

\subsection{Privilege escalation, buffer overflow, rootkits}
\label{sec:privilege-escalation}

U2R attacks start from local access and end with superuser rights. A buffer
overflow, the KDD example, is a programming error that an attacker can
exploit. ATT\&CK files the exploitation of programming errors in a program, a
service or the operating system and its kernel under privilege escalation
\citep{mitre2025t1068}. The match with the buffer overflow is ours. Rootkits
hide programs, files, network connections and other system components by
intercepting and modifying operating system calls, and ATT\&CK files them
under the tactic Stealth \citep{mitre2026t1014}. They act on the host, so the
hiding itself leaves no trace in flow records. This conclusion is ours.

\subsection{Why these classes are the hard ones for data-driven detectors}
\label{sec:r2l-u2r-hard}

Three reasons make R2L and U2R hard for a detector built on flow features.
First, the evidence is in the data portion of a single connection
\citep{kddcup1999task,tavallaee2009kddcup99}. Flow records do not carry it
\citep{sperotto2010flowbased}. The content features designed to capture it
need readable payload, and encryption removes it
\citep{tavallaee2009kddcup99,moriarty2018rfc8404}. Second, these classes are
small in NSL-KDD, so a learner has few examples.
% FIXME(open): R2L and U2R being small classes in NSL-KDD has no source yet; give the class counts in the dataset description (Chapter 9, Sec. 9.2) or cite them.
Third, the logic behind these attacks, a failed login or a privilege change,
concerns a host and not a flow. A rule over flow features can state it only
through indirect signs. The domain knowledge formalised at the end of this
chapter is therefore weaker for these two families than for DoS and Probe, and a
neurosymbolic detector cannot be expected to gain as much on them. This is
our expectation, and the experiments test it.

% ================================================================
% 3.5 Defence mechanisms and where detection sits
% ================================================================
\section{Defence mechanisms and where detection sits}
\label{sec:defences}

\subsection{Scrubbing centres, anycast, rate limiting, BGP flowspec}
\label{sec:mitigation-mechanisms}

A flood that saturates the link of the victim cannot be stopped at the
victim, so the defences act upstream \citep{handley2006rfc4732}. A mitigation
provider, or scrubbing centre, receives the traffic of its customer through a
BGP prefix announcement or through DNS redirection. The cleaned traffic may
come back through a tunnel such as GRE \citep{dobbins2021rfc8903}. Anycast
spreads the load. The same service address is made available at several
nodes, so each node absorbs the attack traffic from its own part of the
network and the work is shared among all nodes \citep{abley2006rfc4786}. Rate
limiting and filtering at the routers can be distributed with BGP Flow
Specification. A rule can match on fields such as the source and destination
prefix, the protocol, the ports, the TCP flags and the packet length, and its
action can limit the rate in bytes or in packets per second
\citep{loibl2020rfc8955}. One condition for a rule to be valid is that its
originator matches the originator of the best-match unicast route for the
destination prefix it names \citep{loibl2020rfc8955}.

Diverting traffic and filtering it both start with the decision that an
attack is under way. Filtering also needs a description of the attack
traffic. The match fields of flow specification are largely the fields of a
flow record. We note this because it is the point where a detector enters.
Since a subtle attack cannot in principle be told from a flash crowd
\citep{handley2006rfc4732}, a rule that is too broad blocks real users.

\subsection{Detection - mitigation - response loop and its timing constraints}
\label{sec:detection-mitigation-loop}

The defences form a loop. In the orchestration use case of DDoS Open Threat
Signaling, telemetry systems alert an orchestrator when they detect an
attack. A mitigation strategy is proposed, in some cases an administrator has
to confirm or select it, and the orchestrator then requests mitigation from
the selected systems \citep{dobbins2021rfc8903}. The GitHub incident shows
the timing. On 28 February 2018 the network monitoring of GitHub detected an
anomaly in the ratio of ingress to egress traffic at 17:21 UTC. GitHub
decided to move its traffic to a mitigation provider, and the command to do
so was initiated at 17:26. Monitoring indicated a full recovery at 17:30. The
site was unavailable from 17:21 to 17:26 and intermittently unavailable from
17:26 to 17:30. GitHub then said it was investigating how to use its
monitoring to enable its mitigation providers automatically
\citep{kottler2018github}. In such a loop the time a detector needs adds
directly to the outage, an argument of ours. The thesis therefore reports
detection latency as a cost. It stops at detection and does not implement the
mitigation, as stated in \cref{sec:scope-limitations}.

% ================================================================
% 3.6 Domain knowledge that can be formalised
% ================================================================
\section{Domain knowledge that can be formalised}
\label{sec:domain-knowledge-formalisation}

\subsection{From informal expert heuristics to first-order logic predicates}
\label{sec:heuristics-to-predicates}

An expert heuristic names a condition on observable quantities and the class
that the condition suggests. Written as logic, the quantities become
predicates over a flow $f$, a window of flows $w$ or a source $s$, and the
heuristic becomes an implication. The predicates stand for features that the
detector computes, such as a packet rate or a count of distinct sources. This
chapter does not fix their thresholds. \Cref{ch:ltn} turns the predicates into
degrees of truth, and \cref{ch:usecases-design} collects the rules in a knowledge base.

\subsection{Table: attack family $\rightarrow$ observable condition $\rightarrow$ candidate axiom}
\label{sec:family-condition-axiom}

\Cref{tab:candidate-axioms} restates the traces of the previous sections as
axioms. They are our formalisation and have not been tested.

\begin{table}[htbp]
\centering
\footnotesize
\setlength{\tabcolsep}{3pt}
\caption{Candidate axioms, our formalisation. DoS is the KDD family and
includes DDoS.}
\label{tab:candidate-axioms}
\begin{tabular}{@{}p{1.9cm}p{3.9cm}p{6.5cm}@{}}
\toprule
Family & Observable condition & Candidate axiom \\
\midrule
UDP/ICMP & High rate, many sources
  & $\mathrm{HighRate}(w)\wedge\mathrm{ManySrc}(w)\rightarrow\mathrm{DoS}(w)$ \\
Reflection & Large replies, service port
  & $\mathrm{ManySrc}(w)\wedge\mathrm{ServicePort}(w)\wedge\mathrm{LargeReplies}(w)\rightarrow\mathrm{DoS}(w)$ \\
SYN flood & SYN excess, incomplete
  & $\mathrm{SynExcess}(w)\wedge\mathrm{Incomplete}(w)\rightarrow\mathrm{DoS}(w)$ \\
Slow DoS & Long flows, low rate
  & $\mathrm{LongFlow}(f)\wedge\mathrm{LowRate}(f)\rightarrow\mathrm{DoS}(f)$ \\
HTTP flood & Complete connections, request rate
  & $\mathrm{Complete}(w)\wedge\mathrm{ReqRate}(w)\rightarrow\mathrm{DoS}(w)$ \\
Probe & Many short flows, many targets
  & $\mathrm{ManyTargets}(s)\wedge\mathrm{Short}(s)\rightarrow\mathrm{Probe}(s)$ \\
R2L & Repeated attempts
  & $\mathrm{RepeatedAttempts}(s)\rightarrow\mathrm{R2L}(s)$ \\
U2R & Nothing in headers & None \\
\bottomrule
\end{tabular}
\end{table}

The table is uneven. The first axiom needs many sources, so it covers DDoS
and not a flood from a single source. R2L covers only password guessing. U2R
has no axiom, for the reasons of \cref{sec:r2l-u2r-hard}, and HTTP/2 rapid
reset has none because its resets are not visible in flow records. An axiom
can also be wrong. The HTTP flood rule fires on a flash crowd
\citep{handley2006rfc4732}. The quality of the encoded knowledge is one of the
conditions studied in RQ2.

% ================================================================
% 3.7 Summary
% ================================================================
\section{Summary}
\label{sec:ch3-summary}

The four KDD families are the labels and ATT\&CK gives them names. Floods
and scans leave rates and counts in flow records, slow and application
attacks leave behaviour per source, and R2L and U2R leave little. Defences
divert, spread or filter traffic. Filtering needs a description of the attack
that a detector must supply in time. Seven candidate axioms follow.